\chapter{Standar Pelayanan Minimal}
\begin{center}
\renewcommand*{\arraystretch}{1.3}
%\newcolumntype{R}[1]{>{\raggedleft\arraybackslash}p{#1}}
\begin{longtable}{rY{20em}Z{6em}Z{6em}Z{6em}Z{8em}Z{7em}}
% prevent longtable show up twice in TOC. set caption for firsthead, then set empty starred
% caption for next heads. in longtable \toprule should be prepended by \\
\caption{Capaian Standar Pelayanan Minimal (SPM) Bidang Kesehatan Kab. Belitung Timur Tahun \tP}
\\ \toprule
\multirow{2}[0]{*}{No} & \multirow{2}[0]{*}{JENIS PELAYANAN{*}} & \multicolumn{3}{c}{MUTU LAYANAN} & \multirow{2}[0]{7em}{\raggedleft MUTU BARANG/ JASA/ SDM (\%)} & \multirow{2}[0]{6em}{\raggedleft MUTU SPM (\%){***}}\\
\cmidrule(l{2pt}r{2pt}){3-5}
   &                    & PEMBILANG & PENYEBUT{**} & MUTU (\%) &                             &                   \\
\midrule
\endfirsthead
\caption*{}
\\ \toprule
\multirow{2}[0]{*}{No} & \multirow{2}[0]{*}{JENIS PELAYANAN{*}} & \multicolumn{3}{c}{MUTU LAYANAN} & \multirow{2}[0]{7em}{\raggedleft MUTU BARANG/ JASA/ SDM (\%)} & \multirow{2}[0]{6em}{\raggedleft MUTU SPM (\%){***}}\\
\cmidrule(l{2pt}r{2pt}){3-5}
&                    & PEMBILANG & PENYEBUT{**} & MUTU (\%) &                             &                   \\
\midrule
\endhead

 1 & Pelayanan kesehatan ibu hamil                          &  1.546 &  1.843 &  83,48 & 100,00 &  87,11 \\
 2 & Pelayanan kesehatan ibu bersalin                       &  1.666 &  1.843 &  90,40 & 100,00 &  92,32 \\
 3 & Pelayanan kesehatan bayi baru lahir                    &  1.646 &  1.843 &  89,31 & 100,00 &  91,45 \\
 4 & Pelayanan kesehatan balita                             &  7.632 &  8.823 &  86,50 & 100,00 &  89,20 \\
 5 & Pelayanan kesehatan pada usia pendidikan dasar         & 18.982 & 19.640 &  96,65 & 100,00 &  97,32 \\
 6 & Pelayanan kesehatan pada usia produktif                & 84.550 & 87.514 &  96,61 & 100,00 &  97,29 \\
 7 & Pelayanan kesehatan pada usia lanjut                   & 11.881 & 14.379 &  82,63 & 100,00 &  86,10 \\
 8 & Pelayanan kesehatan penderita hipertensi               & 20.853 & 24.256 &  85,97 & 100,00 &  88,78 \\
 9 & Pelayanan kesehatan penderita diabetes melitus         &  2.338 &  2.445 &  95,62 & 100,00 &  96,50 \\
10 & Pelayanan kesehatan orang dengan gangguan jiwa berat   &    320 &    320 & 100,00 & 100,00 & 100,00 \\
11 & Pelayanan kesehatan orang terduga tuberkulosis         &  2.482 &  2.262 & 100,00 & 100,00 & 100,00 \\
12 & Pelayanan kesehatan orang dengan risiko terinfeksi HIV &  2.742 &  3.105 &  88,31 & 100,00 &  90,65 \\
\midrule
\multicolumn{2}{r}{Indeks SPM{****}}                        &        &        &        &        & 93,06 \\
                                                     \multicolumn{7}{r}{TUNTAS UTAMA} \\
\bottomrule
\end{longtable}
\par\end{center}


{*}) \emph{Sesuai Peraturan Menteri Kesehatan Nomor 6 Tahun 2024 tentang Standar Teknis Pemenuhan Mutu Pelayanan Dasar Pada Standar Pelayanan
Minimal Bidang Kesehatan}

{**}) \emph{Berdasarkan estimasi dan tidak selalu menggambarkan jumlah yang sebenarnya di populasi}

{***}) \emph{Dihitung dari komponen mutu layanan dan komponen mutu barang/jasa/SDM}

{****}) \emph{Sesuai Peraturan Menteri Dalam Negeri Nomor 59 Tahun 2021 Tentang Penerapan Standar Pelayanan Minimal}
